% Modernised reading — fonds Alexandre Grothendieck, shelfmark 88, pages 1-15.
%
% This is an INTERPRETATION, not a transcription. Everything the critical
% apparatus carried moves into footnotes. In any dispute of substance the
% transcription governs: whoever cites Grothendieck must cite batch-01.fr.tex.
%
% Pass: Opus 5 (claude-opus-5), 2026-09-13 — first pass, unchecked against the
% pages by a human. The folder was read whole; it holds one batch, pages 1-15.
% Page 1 is a wrapper carrying only a title, page 9 its blank other face.
%
% What the folder is. Two runs about maps (« cartes ») and the groups that act
% on their flags. Run 1 (pages 2-8), under the wrapper's title « Cartes
% associées aux syst. de Racines »: a root system is replaced by its Weyl
% group with the torsor of its bases and the simple reflections; the
% parabolic subsets give the Coxeter complex, a triangulated sphere; its
% flags carry operations sigma_0..sigma_{l-1}, and the order of the last
% product is computed to be 2 nu, nu the lcm of the Coxeter exponents. Run 2
% (pages 10-15), « Théorème (à prouver) »: the simply connected quasi-regular
% maps are the standard maps C_{p,q}, their automorphism group is the
% quotient Gamma_{p,q} of the cartographic group, they are realised as
% geodesic tessellations of the sphere, the plane or the hyperbolic plane
% according to (p,q), and every map of type (p,q) is a quotient of C_{p,q}.
%
% Notation fixed for the whole document. Sigma a root system in V, W its Weyl
% group, R the set of its bases (his « repères »), I the type set, phi(r,i)
% the simple reflection of index i for the base r, n_ij the order of
% phi(r,i)phi(r,j), l = Card I. For maps: p the order of the vertices, q the
% order of the faces (he writes C_{p,q}), rho_s and rho_f the rotations about
% a vertex and a face in the cartographic group, Gamma_{p,q} the quotient.
% The Schläfli symbol of C_{p,q} is {q,p}; stated once, not used.
%
% Substantive departures from the pages, each footnoted where it occurs:
% — page 2, d): I with the n_ij is the Coxeter graph, not the Dynkin diagram;
%   B_n and C_n (n >= 3) have the same image, so Phi does not reflect
%   isomorphism, and Aut(Sigma) -> Aut(Phi(Sigma)) is not onto when Sigma has
%   a component B_2, G_2 or F_4, or components B_n and C_n together. Pages 3-4 prove injectivity in general and
%   bijectivity exactly when the two diagrams have the same automorphisms;
%   the top row of page 4's diagram must end in the automorphism group of the
%   Dynkin diagram. The reading states this.
% — page 3: « R partie homogène de u_R » read u_W.
% — page 5: vertices of the Coxeter complex correspond to the MAXIMAL proper
%   parabolic parts; « paraboliques minimaux » is minimal for the reversed
%   order he has just adopted. Footnoted, not a correction.
% — page 6: the first formula defines sigma_{i'} for i' up to l-1 and the
%   second redefines sigma_{l-1}; the reading keeps sigma_i = id x tau_{i+1}
%   for i <= l-2. Page 7: « n ≡ 0 (2 n_..) » read i ≡ 0 (2 n_..).
% — page 14, (c_2): « p + q >= 10 ou (p,q) != (3,6),(6,3) » cannot be right
%   as written; for p, q >= 3 the hyperbolic case is p + q >= 9 with (3,6),
%   (6,3) excluded, i.e. 1/p + 1/q < 1/2.
% — page 15: the underlined proper name read « Bourguignon » is not replaced;
%   the footnote says what the context (maps on compact surfaces, algebraic
%   curves) is known under today, without putting a name in his hand.
%
% Modern names given and footnoted as ours: Coxeter graph, Coxeter complex,
% building of type W, extended triangle group [q,p], regular map and
% Schläfli symbol, uniformisation, Fuchsian group, Belyi's theorem and dessins
% d'enfants. « Groupe cartographique » is his own term (Esquisse d'un
% programme, 1984).
%
% What the folder announces and never establishes: the whole of run 2, which
% he heads « à prouver »; the proof that the flag complex of page 5 is a
% sphere (« manifestement »); b) and the NB of page 13, largely illegible;
% the determination of the pi of page 15.
\documentclass[11pt,a4paper]{article}
% Le préambule commun à toutes les transcriptions du fonds Grothendieck.
%
% Il tient en une idée : **l'incertitude se compose, elle ne se résout pas.**
% Une transcription qui lisse un mot douteux a détruit la seule chose pour
% laquelle on la faisait — un lecteur doit pouvoir savoir, à chaque endroit,
% ce qui était sur le feuillet et ce qui a été ajouté. Les six macros qui
% suivent sont donc l'appareil critique, et non de la décoration.
%
% Les mêmes commandes sont reconnues par `scripts/render.mjs`, qui produit la
% vue de lecture du site. Une macro ajoutée ici doit l'être aussi là-bas,
% faute de quoi le rendu échouera bruyamment — ce qui est voulu : un rendu
% silencieusement faux est pire qu'un rendu absent.

\ProvidesPackage{grothendieck}[2026/08/08 Transcriptions du fonds Grothendieck]

\usepackage{amsmath,amssymb,amsthm}
\usepackage{mathrsfs}
\usepackage{stmaryrd}
% Les diagrammes commutatifs sont partout dans ce fonds : c'est la langue
% même de la géométrie algébrique de Grothendieck, et une transcription qui
% les rendrait en prose ne transcrirait rien.
\usepackage{tikz-cd}
% Français par défaut — c'est la langue des feuillets — mais l'anglais est
% chargé aussi : la traduction et le résumé sont des documents anglais, et la
% typographie française (espaces avant les deux-points) y serait une faute.
% Ces documents basculent par \selectlanguage{english} en tête de corps.
\usepackage[english,french]{babel}
\usepackage{fontspec}
\usepackage{geometry}
\geometry{a4paper,margin=25mm,marginparwidth=28mm}
\usepackage{marginnote}
\usepackage{xcolor}
% ulem, not soul. `soul`'s \st and \ul are built on a character-by-character
% scan that XeLaTeX's font machinery breaks: the document fails with an
% undefined control sequence rather than degrading. ulem does the same two jobs
% and is Unicode-engine safe. `normalem` stops it hijacking \emph, which the
% transcriptions use for Grothendieck's own emphasis.
\usepackage[normalem]{ulem}
\usepackage{hyperref}
\hypersetup{colorlinks=true,linkcolor=black,urlcolor=black,pdfborder={0 0 0}}

% --- Métadonnées du lot -----------------------------------------------------
% Reprises telles quelles par le script de rendu, qui les affiche en tête de
% la vue de lecture. Elles fixent ce que le fichier prétend couvrir : sans
% elles, une transcription flotte sans point d'attache dans un fonds de seize
% mille feuillets.

\newcommand{\folder}[1]{\gdef\grothFolder{#1}}
\newcommand{\batch}[1]{\gdef\grothBatch{#1}}
\newcommand{\leaves}[2]{\gdef\grothFirst{#1}\gdef\grothLast{#2}}
\newcommand{\foldertitle}[1]{\gdef\grothFoldertitle{#1}}
\gdef\grothFolder{?}\gdef\grothBatch{?}\gdef\grothFirst{?}\gdef\grothLast{?}\gdef\grothFoldertitle{}

% --- Le résumé --------------------------------------------------------------
%
% Il ouvre la lecture modernisée, sous le titre « Résumé », et il est composé
% en petit. Ce n'est pas un ornement : c'est le seul passage du document qui
% ne soit pas des mathématiques, et un lecteur doit pouvoir le reconnaître
% avant de l'avoir lu — puis le sauter s'il connaît déjà le sujet, ou s'y
% arrêter s'il ne le connaît pas. Le filet qui le suit marque la reprise du
% corps.

\newenvironment{resume}
  {\small\setlength{\parskip}{0.45em}}
  {\par\medskip\hrule\medskip}

% --- Mots-clés ---------------------------------------------------------------
%
% En anglais, en clôture du résumé de la lecture modernisée : le vocabulaire
% moderne sous lequel chercher ce que les feuillets construisent. Le manifeste
% du site (`npm run manifest`) les extrait pour étiqueter la cote entière —
% cette ligne est la source unique des tags, il n'y a pas de fichier à côté.
\newcommand{\keywords}[1]{\par\smallskip\noindent
  {\footnotesize\textsc{Keywords} --- \textit{#1}}\par}

% --- L'appareil critique ----------------------------------------------------

% Le numéro de feuillet, en marge. C'est l'ancre qui permet de régler un
% désaccord sur une formule en regardant *un* feuillet plutôt qu'une chemise
% de six cents. Le site s'en sert aussi pour tourner les pages du fac-similé
% quand on fait défiler la transcription.
\newcommand{\leaf}[1]{%
  \marginnote{\footnotesize\color{gray!70}\textsf{#1}}%
  \label{leaf:#1}\ignorespaces}

% Illisible : on ne devine pas. Un mot inventé qui se lit comme les autres est
% le pire résultat possible.
\newcommand{\ill}{\textcolor{red!60!black}{[\,\ldots\,]}}

% Lecture douteuse : le mot est donné, et signalé comme incertain.
\newcommand{\uncertain}[1]{\uline{#1}}

% Ajout de l'éditeur — une lettre manquante, un mot sous-entendu.
\newcommand{\add}[1]{\textcolor{blue!50!black}{[#1]}}

% Biffé par Grothendieck lui-même. Ce qu'il a rayé fait partie du document :
% une rature dit souvent où la pensée a changé de direction.
\newcommand{\struck}[1]{\sout{#1}}

% Note du transcripteur, sur l'état matériel du feuillet.
\newcommand{\note}[1]{\par\smallskip{\small\color{gray!60!black}\textit{[#1]}}\par\smallskip}

% Note marginale de Grothendieck, distincte du corps du texte.
\newcommand{\marginal}[1]{\par\smallskip\noindent\hspace{1em}%
  {\small\color{gray!60!black}#1}\par\smallskip}

% --- Titre ------------------------------------------------------------------

\AtBeginDocument{%
  \begin{center}
    {\large\bfseries Fonds Alexandre Grothendieck --- cote n\textsuperscript{o}~\grothFolder}\\[2pt]
    {\itshape\grothFoldertitle}\\[4pt]
    {\small feuillets \grothFirst--\grothLast\ (lot \grothBatch)}\\[2pt]
    {\footnotesize\color{gray!60!black}Transcription établie d'après le fac-similé numérisé
     par l'Université de Montpellier.\\
     Les lectures incertaines sont soulignées, les illisibles notées [\,\ldots\,],
     les ajouts entre crochets.}
  \end{center}
  \vspace{1em}}

\endinput


\folder{88}
\pages{1}{15}
\dating{s.d. — le groupe « Géométrie et topologie combinatoire » (69 à 102) est daté [1976-vers 1986]}
\watermark{Édition de démonstration\\interprétation personnelle de l'œuvre}
\foldertitle{Cartes standards : notes manuscrites (s.d.) — lecture modernisée du dossier entier}

\begin{document}

\section*{Résumé}

\begin{resume}
Une \emph{carte}, au sens de ces feuillets, est un dessin tracé sur une
surface --- des sommets, des arêtes, des faces --- comme la carte d'un pays
découpé en départements. Le cube dessiné sur une sphère en est une ; le pavage
d'un carrelage en est une autre, sur le plan. Grothendieck a proposé de coder
une carte par un objet purement algébrique : l'ensemble de ses \emph{drapeaux}
(un sommet, une arête qui le touche, une face qui la borde) et trois
opérations qui changent l'un des trois éléments en gardant les deux autres. Le
groupe engendré par ces opérations s'appelle le \emph{groupe cartographique},
et toute carte est un ensemble sur lequel il agit.

Le dossier contient deux textes qui tournent autour de cette idée. Le premier
part des systèmes de racines --- les configurations de vecteurs très
symétriques qui classent les groupes de Lie --- et y trouve une carte cachée :
les hyperplans de réflexion découpent la sphère en triangles, et les
opérations sur les drapeaux de ce découpage sont données par les réflexions du
groupe de Weyl. Il calcule l'ordre de la rotation qui en résulte. En chemin il
se demande si le groupe de Weyl, avec ses générateurs, suffit à retrouver le
système de racines ; la réponse qu'il donne est trop optimiste, parce que les
générateurs ne voient pas la différence de longueur entre les racines.

Le second texte est un théorème qu'il s'assigne à démontrer. Parmi toutes les
cartes, les plus symétriques sont celles où tous les sommets ont le même
nombre $p$ d'arêtes et toutes les faces le même nombre $q$ de côtés. Il en
existe exactement une, à la bonne notion d'équivalence près et à quelques
types dégénérés près, par couple $(p, q)$ --- à condition de la dessiner sur une surface simplement connexe ---,
et elle vit sur la sphère, le plan ou le plan hyperbolique selon que
$\frac{1}{p} + \frac{1}{q}$ est plus grand, égal ou plus petit que
$\frac{1}{2}$. Sur la sphère ce sont les cinq solides de Platon et deux
familles dégénérées ; sur le plan, les trois pavages réguliers ; tout le reste
est hyperbolique. Toute autre carte de type $(p,q)$ s'obtient en repliant
celle-là par un groupe.

C'est la géométrie qui, dans l'\emph{Esquisse d'un programme} de 1984, mène aux
« dessins d'enfants » : une carte sur une surface compacte définit une courbe
algébrique, et ces feuillets s'arrêtent au bord de cette observation.

\keywords{root system, Weyl group, Coxeter graph, Coxeter complex, cartographic group, regular map, triangle group, uniformization, Platonic solids}
\end{resume}

\section*{Le fil du dossier, et les conventions}

\begin{enumerate}
\item[1.] \emph{Du système de racines au groupe de Weyl épinglé} (pages 2 à 4) :
un foncteur $\Phi$, et la question de savoir s'il est une équivalence.
\item[2.] \emph{Le complexe des drapeaux} (pages 4 et 5) : parties paraboliques,
facettes, triangulation de la sphère, repères.
\item[3.] \emph{Les opérations sur les repères} (pages 6 à 8) : relations, et
l'ordre $2\nu$.
\item[4.] \emph{Le théorème des cartes standard} (pages 10 à 15), en sept
articles, énoncé et non démontré.
\end{enumerate}

Les deux textes ne se citent pas. Le premier construit, à partir d'un système
de racines de rang $\ell$, une carte de dimension $\ell - 1$ ; le second ne
parle que de cartes de dimension 2. Ce qui les rattache est la méthode ---
lire une carte sur ses drapeaux et les opérations qui les échangent --- et le
titre de la chemise.

\emph{Conventions.} $\Sigma$ est un système de racines dans $V$, $W$ son groupe
de Weyl, $R$ l'ensemble de ses bases --- ses « repères » ---, sur lequel $W$
opère simplement transitivement, $I$ l'ensemble des types, de cardinal
$\ell$, et $\varphi(r, i)$ la réflexion simple d'indice $i$ relative à la base
$r$ ; $n_{ij}$ est l'ordre de $\varphi(r,i)\varphi(r,j)$. Pour une carte, $p$
est l'ordre des sommets et $q$ celui des faces ; le symbole de Schläfli de
$C_{p,q}$ est $\{q, p\}$, et on ne l'utilise pas.\footnote{Il écrit
$\overline{\mathbb{N}}^{*}$ pour $\mathbb{N}^{*} \cup \{\infty\}$, « comb. »,
« gpe », « ppcm », « pt » pour combinatoire, groupe, plus petit commun
multiple, point, et « ½ direct » pour semi-direct. Il surmonte le $\Sigma$
d'une barre, inconstamment, qu'on ne reprend pas. Sur les pages 6 et 7, une
parenthèse ouvrante suivie de son $r$ ressemble à un $w$ ; tous les premiers
arguments de $\varphi$ sont des éléments de $R$.}

\pagerange{2}{4}
\section*{Le groupe de Weyl épinglé (pages 2 à 4)}

À un système de racines $\Sigma$ on associe le quadruplet
\[
  \Phi(\Sigma) = \bigl(W,\ R,\ I,\ \varphi : R \times I \to W\bigr),
\]
qui satisfait :
\begin{enumerate}
\item[a)] $\varphi(wr, i) = \mathrm{int}(w)\,\varphi(r,i)$ ;
\item[b)] pour tout $r$, les $\varphi(r,i)$ engendrent $W$ et en font un
« groupe de Weyl épinglé » ;
\item[c)] $\varphi(r,i)^{2} = 1$ et $\bigl(\varphi(r,i)\varphi(r,j)\bigr)^{n_{ij}}
= 1$ pour $i \neq j$ forment une présentation de $W$ ;
\item[d)] $I$, muni des $n_{ij}$, $i$ étant relié à $j$ si et seulement si
$n_{ij} \geq 3$, est le graphe de Coxeter d'un groupe de Weyl.
\end{enumerate}
Les systèmes de ce type forment une catégorie, et $\Sigma \mapsto \Phi(\Sigma)$
est un foncteur.\footnote{La page appelle « diagramme de Dynkin » le graphe de
d). Ce n'en est pas un : le diagramme de Dynkin porte en outre, sur les liaisons
doubles et triples, l'orientation qui distingue les racines longues des
courtes, et les $n_{ij}$ n'en gardent rien. C'est ce qu'on appelle aujourd'hui
le graphe de Coxeter. La distinction change la réponse à la question de la page,
comme on va le voir.}

\emph{Le centre.} Chaque base $r$ donne un système de générateurs
$\Psi(r) = (\varphi(r,i))_{i \in I}$ de $W$ indexé par $I$. On a
$\Psi(wr) = \Psi(r)$ si et seulement si $\mathrm{int}(w)$ fixe tous les
$\varphi(r,i)$, c'est-à-dire si et seulement si $w$ est central. Donc $\Psi$
est injective --- et $R$ s'identifie à un ensemble de systèmes de générateurs
de $W$ indexés par $I$ --- si et seulement si $Z(W) = \{1\}$.\footnote{Un bloc
d'une dizaine de lignes, annulé, précède cette conclusion ; il cherchait le
stabilisateur d'un repère. Le centre de $W$ n'est pas trivial, par exemple,
dès que $-1 \in W$ : types $A_{1}$, $B_{n}$, $D_{2k}$, $E_{7}$, $E_{8}$, $F_{4}$,
$G_{2}$.}

\emph{La question.} $\Phi$ est-il une équivalence de catégories ? La page
répond : il est évidemment fidèle et essentiellement surjectif ; deux systèmes
dont les images sont isomorphes ont des diagrammes isomorphes, donc sont
isomorphes ; il reste la pleine fidélité, c'est-à-dire que
$\mathrm{Aut}(\Sigma) \to \mathrm{Aut}(\Phi(\Sigma))$ soit bijectif.

\emph{Ce qui est démontré} (pages 3 et 4). Soit
$\alpha : W \to \mathrm{Aut}(\Phi(\Sigma))$, $\alpha(w) = (\mathrm{int}(w),
w_{R}, \mathrm{id}_{I})$ ; c'est un morphisme, par a). Il est injectif, et son
image est formée des automorphismes $(u_{W}, u_{R}, \mathrm{id}_{I})$ qui
induisent l'identité sur $I$. En effet, un tel automorphisme est la donnée d'un
automorphisme $u_{W}$ de $W$ et d'une bijection $u_{R}$ de $R$ affine de partie
homogène $u_{W}$,\footnote{La page écrit « $R$ “partie homogène” de $u_{R}$ » ;
le sens demande $u_{W}$.} qui vérifient
\[
  (\ast\ast) \qquad \varphi(u_{R}\,r,\ i) = u_{W}\,\varphi(r,i).
\]
Choisissons $r_{0} \in R$, et $a \in W$ tel que $u_{R}(r_{0}) = a\,r_{0}$ ;
alors $u_{R}(w\,r_{0}) = u_{W}(w)\,a\,r_{0}$, et $(\ast\ast)$ pour
$r = w\,r_{0}$ se développe en
\[
  a\,\varphi(r_{0}, i)\,a^{-1} = u_{W}\bigl(\varphi(r_{0}, i)\bigr)
  \qquad (i \in I).
\]
Les $\varphi(r_{0}, i)$ engendrant $W$, $u_{W} = \mathrm{int}(a)$ ; puis
$u_{R}(w\,r_{0}) = (awa^{-1})(ar_{0}) = a(w\,r_{0})$, soit $u_{R} = a_{R}$, et
l'automorphisme est $\alpha(a)$.

On a donc un diagramme à lignes exactes

\begin{tikzcd}[column sep=small, row sep=small, nodes={font=\scriptsize}]
  1 \arrow[r] & W \arrow[r, "\alpha_0"] \arrow[d, "\mathrm{id}"] &
    \operatorname{Aut}(\Sigma) \arrow[r, "p_0"] \arrow[d] &
    \operatorname{Aut}(\mathrm{Dyn}(\Sigma)) \arrow[r] \arrow[d, hook] & 1 \\
  1 \arrow[r] & W \arrow[r, "\alpha"'] &
    \operatorname{Aut}(\Phi(\Sigma)) \arrow[r, "p"'] &
    \operatorname{Aut}(I, (n_{ij})) &
\end{tikzcd}

où $\mathrm{Aut}(\Sigma)/W = \mathrm{Aut}(\mathrm{Dyn}(\Sigma))$ est le groupe
des automorphismes du diagramme de Dynkin --- « gr.\ des automorphismes ext.\
de $\Sigma$ », écrit-il --- et $\mathrm{Aut}(I, (n_{ij}))$ celui du graphe de
Coxeter. Le lemme des cinq donne :

\textbf{Proposition.} $\mathrm{Aut}(\Sigma) \to \mathrm{Aut}(\Phi(\Sigma))$ est
injectif ; il est bijectif si et seulement si tout automorphisme du graphe de
Coxeter de $\Sigma$ respecte les longueurs des racines, c'est-à-dire si et
seulement si aucune composante irréductible de $\Sigma$ n'est de type $B_{2}$,
$G_{2}$ ou $F_{4}$, et si $\Sigma$ n'a pas à la fois une composante de type
$B_{n}$ et une de type $C_{n}$, $n \geq 3$.

\textbf{Conséquence.} $\Phi$ est fidèle, et essentiellement surjectif sur les
systèmes qui satisfont a) à d). Il n'est pas une équivalence : il n'est pas
plein dans ces cas, et il ne distingue pas $B_{n}$
de $C_{n}$ pour $n \geq 3$. Sa restriction aux systèmes de racines dont toutes
les racines ont même longueur --- types $A$, $D$, $E$ --- est une équivalence
sur son image.\footnote{La page écrit le même $\mathrm{Aut}(I)$ au bout des
deux lignes, avec deux flèches identiques, et conclut « le diagramme des 5 […]
montre alors que la flèche verticale médiane est un iso, cqfd ». Cela suppose
que les automorphismes de $I$ qui respectent les $n_{ij}$ proviennent tous
d'automorphismes de $\Sigma$, ce qui est faux pour $B_{2}$ (les deux sommets
échangés, $n_{12} = 4$), $G_{2}$ ($n_{12} = 6$) et $F_{4}$ (le retournement du
graphe), où la symétrie du graphe de Coxeter échange racines longues et
courtes. De même, « des diagrammes isomorphes, donc sont isomorphes » n'est
vrai que pour les diagrammes de Dynkin, et $W(B_{n}) = W(C_{n})$ avec la même
présentation. Pour ce qui suit dans le dossier --- des cartes, qui ne
dépendent que de $W$ et de ses générateurs --- la perte est sans
conséquence.}

\pagerange{4}{5}
\section*{Le complexe des drapeaux (pages 4 et 5)}

Les parties paraboliques de $\Sigma \subset V$ correspondent aux facettes
fermées de l'arrangement des hyperplans de réflexion dans le dual
$\check{V}$, et l'inclusion des parties paraboliques --- ou des sous-groupes
paraboliques, quand $\Sigma$ provient d'un groupe algébrique --- devient
l'inclusion \emph{opposée} des facettes. On munit l'ensemble des parties
paraboliques propres de cet ordre opposé ; il devient la géométrie des
\emph{drapeaux stricts}, un complexe simplicial dont les sommets sont les
paraboliques minimaux pour cet ordre, et les simplexes maximaux les
« Borels », ou systèmes de racines simples, ou repères
$r \in R$.\footnote{Minimaux pour l'ordre opposé, c'est-à-dire maximaux pour
l'inclusion : les sommets sont les parties paraboliques propres maximales,
qui correspondent aux « drapeaux minimaux », comme il l'ajoute en interligne.
Le nom du complexe, après « correspondant à une un », ne se lit pas ; la même
lettre $V$ est écrite là où les sommets et les simplexes maximaux sont
nommés.}

« Il s'agit manifestement d'une triangulation de la sphère $S^{\ell - 1}$ » :
les chambres de Weyl découpent la sphère unité de $\check{V}$ en simplexes
sphériques, à supposer que $\Sigma$ engendre $V$, de dimension
$\ell$.\footnote{C'est le complexe de Coxeter de $W$ --- l'immeuble de type $W$
de rang $\ell$ réduit à un seul appartement ; les noms ne sont pas sur la
page.}

\emph{Les repères du complexe.} Un drapeau complet de faces d'un simplexe
maximal $r$ est une numérotation de ses $\ell$ sommets, et les sommets d'une
chambre sont indexés par leurs types. L'ensemble des repères est donc
\[
  R \times \mathrm{Rep}'(I),
  \qquad
  \mathrm{Rep}'(I) = \mathrm{Bij}\bigl([1, \ell],\ I\bigr),
\]
et $\mathrm{Rep}'(I)$ s'envoie sur l'ensemble des ordres totaux sur $I$.
C'est un torseur à droite sous $\mathfrak{S}_{\ell}$.

\pagerange{6}{8}
\section*{Les opérations sur les repères, et l'ordre $2\nu$ (pages 6 à 8)}

Sur les repères d'un complexe de dimension $\ell - 1$ opèrent $\ell$
involutions $\sigma_{0}, \ldots, \sigma_{\ell-1}$, $\sigma_{k}$ changeant la
face de dimension $k$ du drapeau en gardant les autres. Ici :
\[
  \sigma_{i}(r, u) = (r,\ u \circ \tau_{i+1}) \quad (0 \leq i \leq \ell - 2),
  \qquad
  \sigma_{\ell-1}(r, u) = \bigl(\varphi(r, u(\ell)) \cdot r,\ u\bigr),
\]
où $\tau_{i+1} \in \mathfrak{S}_{\ell}$ est la transposition de $i+1$ et
$i+2$.\footnote{La page définit d'abord
$\sigma_{i'}(r, u) = (r, u \circ \tau_{i'})$ pour $0 \leq i' \leq \ell - 1$,
puis redéfinit $\sigma_{\ell-1}$ à la ligne suivante ; elle annonce aussi
« $\sigma_{0}, \ldots, \sigma_{\ell}$ ». On retient le décalage d'indice de sa
phrase précédente, $\sigma_{i} = \mathrm{id}_{R} \times \tau_{i+1}$, qui est
celui qu'utilisent tous les calculs.} La seconde formule traduit la
géométrie : la cloison du drapeau est la face de $r$ opposée au sommet de type
$u(\ell)$, et la chambre voisine à travers elle est $s_{u(\ell)}\,r$.

\emph{Relations.} On vérifie :
\begin{itemize}
\item $\sigma_{i}^{2} = 1$ pour tout $i$ --- pour $\sigma_{\ell-1}$, parce que
$\varphi(sr, u(\ell)) = s\,\varphi(r,u(\ell))\,s^{-1} = s$ si
$s = \varphi(r, u(\ell))$ ;
\item $\sigma_{i}\sigma_{j} = \sigma_{j}\sigma_{i}$ pour $j \geq i + 2$, le seul
cas non trivial étant $j = \ell - 1$, $i \leq \ell - 3$, où
$(u \circ \tau_{i+1})(\ell) = u(\ell)$ ;
\item $(\sigma_{i}\sigma_{i+1})^{3} = 1$ pour $i \leq \ell - 3$, par
$(\tau_{i+1}\tau_{i+2})^{3} = 1$.
\end{itemize}

\emph{Le dernier produit.} Soit $\pi = \sigma_{\ell-1}\sigma_{\ell-2}$ :
\[
  \pi(r, u) = \bigl(\varphi(r, u(\ell - 1)) \cdot r,\ u \circ \tau_{\ell-1}\bigr).
\]
Posons $a = \varphi(r, u(\ell-1))$ et $b = \varphi(r, u(\ell))$. Par a),
$\pi^{2}(r,u) = \bigl(a\,\varphi(r,u(\ell))\,a^{-1}\,a\,r,\ u\bigr) = (ab\,r, u)$,
et par récurrence
\[
  \pi^{2n}(r, u) = \bigl((ab)^{n}\, r,\ u\bigr),
  \qquad
  \pi^{2n+1}(r, u) = \bigl((ab)^{n} a\, r,\ u \circ \tau_{\ell-1}\bigr).
\]
Une puissance impaire change $u$ ; une puissance paire $2n$ fixe $(r, u)$ si et
seulement si $(ab)^{n} = 1$, $R$ étant un torseur, c'est-à-dire si et seulement
si $n_{u(\ell-1),\,u(\ell)}$ divise $n$. Sur l'orbite de $(r,u)$, $\pi$ est donc
d'ordre $2\,n_{u(\ell-1),\,u(\ell)}$.\footnote{La page écrit la conclusion
« $n \equiv 0\ (2\,n_{u(\ell-1),\,u(\ell)})$ » ; c'est $i = 2n$ qui est
multiple de $2\,n_{u(\ell-1),u(\ell)}$. Géométriquement, une face de
codimension 2 de type $I - \{i,j\}$ est contenue dans $2n_{ij}$ chambres.}
Faisant varier $u$, toutes les paires $i \neq j$ apparaissent, et :

\textbf{Proposition.} L'ordre de $\sigma_{\ell-2}\sigma_{\ell-1}$ sur
$R \times \mathrm{Rep}'(I)$ est $2\nu$, où $\nu$ est le plus petit commun
multiple des $n_{ij}$, $i \neq j$.\footnote{La page 8 écrit « ordre de
$\pi = \sigma_{\ell-2}\sigma_{\ell-1} = 2\nu$ », égalant l'ordre à $\pi$ ; et
$\sigma_{\ell-2}\sigma_{\ell-1}$ est l'inverse du $\pi$ de la page 6, ce qui
ne change pas l'ordre. Les $n_{ij} = 2$ comptent dans le ppcm.}

Pour $\ell = 3$ c'est une carte de dimension 2 sur la sphère : les faces sont
des triangles ($\sigma_{0}\sigma_{1}$ d'ordre 3), et un sommet de type $k$ est
d'ordre $2 n_{ij}$, $\{i, j\} = I - \{k\}$ ; elle n'est quasi-régulière, au
sens de la seconde partie, que si les trois $n_{ij}$ sont égaux.

\pagerange{10}{10}
\section*{Le théorème des cartes standard : (1) à (4) (page 10)}

La seconde partie porte, de sa main, le titre « Théorème (à prouver) ». Une
\emph{carte cellulaire} est un triple $(X \supset K \supset S)$ : une surface,
un graphe tracé dessus dont le complémentaire est une réunion de cellules, et
l'ensemble de ses sommets. Elle est \emph{quasi-régulière} si tous ses sommets
combinatoires ont le même ordre $p$ et toutes ses faces le même ordre $q$, de
\emph{type} $(p, q)$, $1 \leq p, q \leq \infty$ ; \emph{régulière} si son groupe
d'automorphismes est transitif sur les repères ; \emph{complète} si les sommets
d'ordre fini sont effectifs.

Dans le langage de l'\emph{Esquisse d'un programme}, les repères --- ou
bi-drapeaux --- d'une carte sont permutés par le \emph{groupe cartographique}
$\widehat{\Gamma} = \langle \sigma_{0}, \sigma_{1}, \sigma_{2} \mid
\sigma_{i}^{2},\ (\sigma_{0}\sigma_{2})^{2} \rangle$, et
$\rho_{s} = \sigma_{1}\sigma_{2}$, $\rho_{f} = \sigma_{0}\sigma_{1}$ sont les
rotations autour d'un sommet et d'une face.\footnote{Les notations
$\widehat{\Gamma}$, $\rho_{s}$, $\rho_{f}$ et le mot « groupe cartographique »
sont sur la page ; la présentation est celle de l'\emph{Esquisse}, qui n'y est
pas écrite. L'ordre des indices $\sigma_{i}$ est une convention que les
feuillets ne fixent pas.}

\begin{enumerate}
\item[(1)] Toute carte cellulaire non vide, connexe, simplement connexe et
quasi-régulière est régulière ; elle est complète, sauf si elle est isomorphe
à $(\mathbb{R}^{2} \supset \mathbb{R})$, dont la complétion est la sphère munie
d'un grand cercle ponctué.
\item[(2)] Une telle carte régulière et complète --- une \emph{carte
standard} --- est déterminée à isomorphisme près par son type $(p, q)$, et
tous les types sont atteints sauf $(1, p)$ et $(p, 1)$ pour $p \neq 2$. On
note $(C_{p,q}, r_{p,q})$ la carte standard épinglée par un repère, définie à
isomorphisme unique près dans la catégorie des cartes à isotopie
près.\footnote{« Isotopique » est lu avec doute. La page écrit « correspondant à
$(q,p)$ » avec les indices dans l'autre ordre, et biffe « homéomorphisme près »
pour « isomorphisme près ».}
\item[(3)] $\mathrm{Aut}(C_{p,q}) \simeq \widehat{\Gamma}/\langle \rho_{s}^{p},
\rho_{f}^{q} \rangle = \Gamma_{p,q}$, le sous-groupe distingué engendré par
$\rho_{s}^{p}$ et $\rho_{f}^{q}$ étant sous-entendu, et aucune relation n'étant
imposée pour un ordre infini.\footnote{$\Gamma_{p,q}$ est le groupe de
réflexions triangulaire étendu $[q, p]$, engendré par trois réflexions dont
les produits deux à deux sont d'ordres $2$, $p$ et $q$ ; le nom n'est pas sur la
page. Un second numéral cerclé, entre (3) et (4), est noirci.}
\item[(4)] \textbf{Corollaire.} Pour une carte cellulaire complète connexe
quasi-régulière $\mathcal{C}$ de type $(p,q)$, le groupe fondamental de sa
surface s'identifie à un sous-groupe de $\Gamma_{p,q}$ ; si $\mathcal{C}$ est
régulière et épinglée, le groupe de ses automorphismes s'identifie au quotient
correspondant de $\Gamma_{p,q}$.\footnote{L'énoncé court sur deux pages et
plusieurs de ses mots sont illisibles ; la seconde moitié en particulier est
reconstituée à partir de « régulière épinglée », « automorphismes », « groupe
fondamental » et « $\Gamma_{p,q}$ ». La forme précise en est donnée par (7).}
\end{enumerate}

\pagerange{11}{13}
\section*{(5) : surfaces standard et cartes géodésiques (pages 11 à 13)}

\emph{Surfaces standard.} Une \emph{surface riemannienne standard} est une
surface analytique réelle munie d'une métrique riemannienne, non vide,
connexe, simplement connexe, dont le groupe des isométries est transitif. À un
facteur constant près sur la métrique, c'est l'une des trois :

\begin{enumerate}
\item[1°)] la sphère $S^{2}$, espace homogène sous $SO(3, \mathbb{R})$ ---
plan sphérique, déplacements sphériques ;
\item[2°)] le plan complexe, sous le produit semi-direct
$\mathrm{U} \cdot \mathbb{C}$ du cercle unité et des translations, groupe des
déplacements euclidiens --- plan euclidien ;
\item[3°)] le demi-plan de Poincaré, espace homogène sous
$SL(2, \mathbb{R})/\{\pm 1\}$ --- plan hyperbolique.
\end{enumerate}

Les groupes écrits sont les composantes neutres des groupes d'isométries,
c'est-à-dire les isométries qui conservent l'orientation, ou les
automorphismes de la structure complexe associée à une
orientation.\footnote{Les trois qualificatifs --- sphérique, euclidien,
hyperbolique --- sont écrits en biais dans la marge. La page qualifie la
métrique d'un « $\geq 0$ » qu'on ne reprend pas : une métrique riemannienne est
définie positive. La classification suit du théorème d'uniformisation, ou
directement de ce qu'une surface homogène est à courbure constante ; le nom
n'est pas sur la page.} Dans les cas 1° et 3° toute similitude est une
isométrie, et l'on normalise la courbure à $\pm 1$ ; dans le cas euclidien
on normalisera la carte elle-même.

\emph{Cartes géodésiques.} Une carte sur une surface riemannienne est
\emph{géodésique} si ses arêtes sont des arcs de géodésiques ;
\emph{métriquement régulière} si elle est géodésique et si le groupe de ses
automorphismes isométriques est transitif sur les repères ; une \emph{carte
géodésique standard} si elle est de plus complète, sur une surface standard.
On doit avoir :

\begin{enumerate}
\item[a)] toute carte standard est topologiquement isomorphe à une carte
géodésique standard $(X_{0}, K_{0}, S_{0})$ ;
\item[b)] entre deux cartes géodésiques standard, il y a dans chaque classe
d'isotopie d'isomorphismes topologiques un unique isomorphisme qui soit une
isométrie --- à un facteur près dans le cas euclidien ;
\item[c)] le groupe des automorphismes d'une carte géodésique standard opère
proprement sur la surface.
\end{enumerate}

Une carte géodésique standard est \emph{normalisée} si, dans le cas euclidien,
ses arêtes de longueur finie sont de longueur 1, les deux autres cas étant
normalisés par la courbure.\footnote{Le NB de la page 13 est écrit sur trois
niveaux d'insertions ; seuls « normalisée », « euclidien », « longueurs des
arêtes » et « égale à 1 » s'y lisent. La restriction de b) au cas non
sphérique, que la page semble faire, n'est pas lisible.}

\textbf{Corollaire.} Le foncteur « métrique » de la catégorie des cartes
géodésiques standard normalisées --- les morphismes étant les isométries qui
respectent la carte --- dans la catégorie des cartes cellulaires standard est
une équivalence de catégories. Épinglées, les cartes géodésiques standard
normalisées vivent canoniquement sur $S^{2}$, $\mathbb{R}^{2} = \mathbb{C}$ et
$\mathbb{H}$ munis de leur repère habituel, et on les note encore
$C_{p,q}$.

\pagerange{14}{14}
\section*{(6) : sphérique, euclidien, hyperbolique (page 14)}

\begin{enumerate}
\item[(6)] $C_{p,q}$ est \emph{sphérique} exactement pour $p = 2$,
$1 \leq q < \infty$ ; $q = 2$, $1 \leq p < \infty$ ; et
\[
  (3,3), \quad (3,4), \quad (4,3), \quad (3,5), \quad (5,3),
\]
le tétraèdre, le cube, l'octaèdre, le dodécaèdre et l'icosaèdre ; ou encore si
et seulement si $C_{p,q}$ est finie, si et seulement si $\Gamma_{p,q}$ est fini.
\item[(b)] $C_{p,q}$ est \emph{euclidienne} exactement pour
\[
  (4,4), \quad (3,6), \quad (6,3),
\]
les pavages carré, hexagonal et triangulaire ; $\Gamma_{p,q}$ est alors infini
et résoluble.
\item[(c)] $C_{p,q}$ est \emph{hyperbolique} dans les autres cas : pour
$p, q \geq 3$, dès que $p + q \geq 9$ et $(p,q) \neq (3,6), (6,3)$, ou dès que
$p$ ou $q$ est infini ; ou encore si et seulement si $\Gamma_{p,q}$ contient un
sous-groupe libre à deux générateurs, donc des sous-groupes libres à $n$
générateurs pour tout $n$.
\end{enumerate}

Les trois cas sont ceux où $\frac{1}{p} + \frac{1}{q}$ est supérieur, égal ou
inférieur à $\frac{1}{2}$ : c'est la somme des angles du triangle
fondamental, $\frac{\pi}{2} + \frac{\pi}{p} + \frac{\pi}{q}$, comparée à
$\pi$.\footnote{Le critère angulaire est le nôtre ; la page donne les listes.
Elle donne aussi deux descriptions du cas hyperbolique, $(c_{1})$ « $(p,q) >
(4,4)$ ; ou $p = 3$, $q \geq 7$ ; ou $q = 3$, $p \geq 7$ ; ou $(4,5)$, $(5,4)$ »,
qui est juste, et $(c_{2})$ « $p + q \geq 10$ ou $(p,q) \neq (3,6), (6,3)$ »,
qui ne peut l'être telle qu'elle est écrite --- la seconde condition, jointe
par « ou », engloberait les cas sphériques ; on la remplace par la forme
équivalente à $(c_{1})$. Les types $(2, \infty)$ et $(\infty, 2)$, pour
lesquels $\frac{1}{p} + \frac{1}{q} = \frac{1}{2}$, sont écrits puis biffés dans
la liste euclidienne ; le premier est la carte $(\mathbb{R}^{2} \supset
\mathbb{R})$ de l'exception de (1), et la page ne dit pas où elle les range.
Le nom du pavage hexagonal est illisible. Que $\Gamma_{p,q}$ contienne un groupe
libre non abélien dans le cas hyperbolique tient à ce qu'il y est un groupe
fuchsien non élémentaire.}

\pagerange{14}{15}
\section*{(7) : les cartes de type $(p,q)$ comme quotients (pages 14 et 15)}

\begin{enumerate}
\item[(7)] Les cartes quasi-régulières épinglées de type $(p,q)$ s'obtiennent
canoniquement à partir des sous-groupes $\pi \subset \Gamma_{p,q}$ opérant
librement sur $X_{p,q}$, la carte étant le quotient de $C_{p,q}$ par $\pi$ et
le groupe fondamental de sa surface s'identifiant à $\pi$ ; la surface
quotient hérite d'une métrique riemannienne pour laquelle la carte est
géodésique, et la carte est régulière exactement quand $\pi$ est distingué
dans $\Gamma_{p,q}$.\footnote{La page écrit « les cartes régulières », puis
« $\Gamma_{p,q}$ est invariant » ; on lit la seconde expression comme la
condition de régularité, ce qui oblige à lire la première comme
« quasi-régulières ». Une partie de l'énoncé est illisible.}
\end{enumerate}

\textbf{NB.} Pour $(p,q)$ donné, déterminer les $\pi$ possibles revient à
étudier l'action de $\Gamma_{p,q}$ sur $X_{p,q}$, et en particulier les
stabilisateurs des points : $\pi$ opère librement si et seulement s'il ne
rencontre ces stabilisateurs qu'en l'élément neutre. Ce sont des groupes
finis : diédraux d'ordre $2p$, $4$ et $2q$ aux sommets, aux milieux d'arêtes et
aux centres de faces, d'ordre 2 aux autres points des droites de
réflexion.\footnote{La description des stabilisateurs est la nôtre ; la page
dit seulement qu'il faut étudier « les stabilisateurs des $\Gamma_{p,q}$ des
pts de $X_{p,q}$ ». Pour un ordre infini le stabilisateur correspondant est
vide de sens, le sommet ou le centre de face étant rejeté à l'infini.}

\emph{Le cas compact.} La page se termine sur un palimpseste dont on lit
qu'une carte finie orientée, $\pi$ étant d'indice fini, définit sur sa surface
une structure complexe, et que dans le cas compact on obtient ainsi des
courbes algébriques projectives. Le dernier mot, souligné, est un nom propre
lu « Bourguignon ».\footnote{Le tiers gauche de la page, écrit en biais, n'est
pas recouvrable, et la prose de la colonne de droite ne l'est qu'en partie.
L'observation qu'on lit est le point de départ de ce que l'\emph{Esquisse
d'un programme} (1984) appelle les dessins d'enfants : une carte sur une
surface compacte orientée définit une courbe algébrique, et, par le théorème de
Belyi (1979), les courbes ainsi obtenues sont exactement celles qui sont
définies sur un corps de nombres. Rien sur la page ne permet de lire le nom de
Belyi à la place de celui qui est écrit, et on ne le substitue pas.}

\end{document}
